\par\addvspace{12pt}\Needspace{9\baselineskip}

\section{方差}

% ITEM: probability-example-019
\solutionitem{probability-example-019}{例6.19}
\(X\)以相同概率取\(-1,1\)，\(Y\)以相同概率取\(-3,3\)。比较它们的期望、方差与标准差，再比较\(2X+5\)和\(Y+5\)。
\solutionanswer
\(X,Y\)期望均为零，方差分别\(1,9\)，标准差\(1,3\)。\(2X+5,Y+5\)期望均为\(5\)，方差分别\(4,9\)，标准差\(2,3\)。平移改变平均位置而不改变波动，数乘则按平方改变方差。


% ITEM: probability-example-020
\solutionitem{probability-example-020}{例6.20}
（据2019年浙江卷改编）
设\(0<a<1\)，\(X\)等可能取\(0,a,1\)。研究\(a\)增大时\(\operatorname{Var}X\)怎样变化，并求其最小值。
\solutionanswer
\(EX=(a+1)/3\)、\(E(X^2)=(a^2+1)/3\)，故\(\operatorname{Var}X=2(a^2-a+1)/9=2(a-1/2)^2/9+1/6\)。随\(a\)在\((0,1/2)\)上增大而减小，在\((1/2,1)\)上增大而增大，最小值\(1/6\)在\(a=1/2\)取得。


% ITEM: probability-example-021
\solutionitem{probability-example-021}{例6.21}
两种测量方法的误差为\(X,Y\)，均值都为0，方差分别为1、4，协方差为0。将结果按权重\(a,1-a\)平均，误差为\(aX+(1-a)Y\)。求使该误差方差最小的权重，并与等权平均比较。
\solutionanswer
方差为\(a^2+4(1-a)^2=5(a-4/5)^2+4/5\)，最优权重\(a=4/5\)，即给低方差方法较高权重，最小方差\(4/5\)。等权平均的方差\(5/4\)，明显更大。这里只需协方差为零，不必另加独立假设。


% ITEM: probability-example-022
\solutionitem{probability-example-022}{例6.22}
\(X\)在\(-2,0,2\)处的概率分别为\(1/8,3/4,1/8\)。分别求\(P(|X|\ge1)\)、\(P(|X|\ge2)\)的实际值与Chebyshev上界。比较两种情形，判断何时界能取到等号。
\solutionanswer
均值零、方差\(4(1/8+1/8)=1\)。两个实际尾概率均为\(1/4\)，Chebyshev上界分别为\(1\)、\(1/4\)。阈值\(2\)时恰取等号，因为非零偏差的平方均为\(4\)，被截去的部分平方偏差也恰为零；阈值\(1\)时界较宽。


% ITEM: probability-exercise-11-01
\solutionitem{probability-exercise-11-01}{习题6.11（1）}
\(X\)以\(1/4,1/2,1/4\)的概率取0、1、2。求\(EX\)、\(\operatorname{Var}X\)，并求使\(E(X-c)^2\)最小的\(c\)。
\solutionanswer
\(EX=1,E(X^2)=3/2\)，方差\(1/2\)。由\(E(X-c)^2=\operatorname{Var}X+(EX-c)^2\)，唯一最优\(c=1\)，最小均方误差\(1/2\)。


% ITEM: probability-exercise-11-02
\solutionitem{probability-exercise-11-02}{习题6.11（2）}
设\(X,Y\)独立，期望分别为2、3，方差分别为1、4。求\(E(2X-Y)\)、\(\operatorname{Var}(2X-Y)\)及\(E(XY)\)。
\solutionanswer
期望\(2\times2-3=1\)，方差\(4\times1+4=8\)，乘积期望\(EX\,EY=6\)。后两式的零协方差和乘积期望分解均由独立性保证。


% ITEM: probability-exercise-11-03
\solutionitem{probability-exercise-11-03}{习题6.11（3）}
设\(\operatorname{Var}X=4\)、\(\operatorname{Var}Y=9\)、\(\operatorname{Cov}(X,Y)=-3\)。求相关系数及\(\operatorname{Var}(X+Y)\)。
\solutionanswer
相关系数\(-3/(2\times3)=-1/2\)，和的方差\(4+9+2(-3)=7\)。


% ITEM: probability-exercise-11-04
\solutionitem{probability-exercise-11-04}{习题6.11（4）}
某非负变量的期望为3，求\(P(X\ge12)\)的上界。若还知道方差为1，使用Chebyshev不等式重新估计并比较。
\solutionanswer
由非负变量的Markov界，\(P(X\ge12)\le3/12=1/4\)。若方差为1，则\(X\ge12\)蕴含\(|X-3|\ge9\)，Chebyshev给上界\(1/81\)，更紧。


% ITEM: probability-exercise-11-05
\solutionitem{probability-exercise-11-05}{习题6.11（5）}
构造两个均值均为0、方差均为1的变量，使其和的方差分别为0、2、4。说明各构造中的相关关系。
\solutionanswer
取等可能的\(X=\pm1\)。令\(Y=-X\)时和方差零，完全负相关；令\(Y\)为另一个独立的等可能\(\pm1\)变量时和方差\(2\)，相关系数零；令\(Y=X\)时和方差\(4\)，完全正相关。


% ITEM: probability-exercise-11-06
\solutionitem{probability-exercise-11-06}{习题6.11（6）}
\(n\)次独立测量的误差均值为0、方差为9。求误差总和与平均误差的方差，并用Chebyshev界确定使平均误差绝对值至少1的概率不超过\(0.05\)的充分测量次数。
\solutionanswer
独立性使总误差方差为\(9n\)，平均误差方差为\(9/n\)。Chebyshev给\(P(|\overline X|\ge1)\le9/n\)，取\(n\ge180\)可保证不超过\(0.05\)。


% ITEM: probability-example-023
\solutionitem{probability-example-023}{例6.23 *}
四个样本点的概率依次为\(1/10,2/10,3/10,4/10\)，\(X\)的值依次为1、3、4、8。只记录样本点属于前两个点还是后两个点，求条件平均组成的变量\(M\)，并验证残差与两个组的指示变量都正交。
\solutionanswer
两组平均分别为\((1+2\times3)/3=7/3\)、\((3\times4+4\times8)/7=44/7\)，所以\(M\)依次取\(7/3,7/3,44/7,44/7\)。残差为\((-4/3,2/3,-16/7,12/7)\)。其与第一组指示量的内积为\((1/10)(-4/3)+(2/10)(2/3)=0\)，与第二组的内积为\((3/10)(-16/7)+(4/10)(12/7)=0\)，验证了正交性。


% ITEM: probability-exercise-12-01
\solutionitem{probability-exercise-12-01}{习题6.12（1） *}
四个等可能样本点上，\(X\)的值依次为1、3、4、8。按前两个点与后两个点分组，求\(E(X\mid\mathcal P)\)，并核对总方差分解。
\solutionanswer
条件平均\(M=(2,2,6,6)\)。\(EX=4\)，\(\operatorname{Var}X=(1+9+16+64)/4-16=13/2\)。两组内方差分别为\(1,4\)，平均为\(5/2\)；组均值方差\(\operatorname{Var}M=4\)，所以\(13/2=5/2+4\)，与总方差分解一致。


% ITEM: probability-exercise-12-02
\solutionitem{probability-exercise-12-02}{习题6.12（2） *}
在（1）的模型中，把后两个点进一步分别记录，求新增信息使最小平方误差减少多少，并与两种条件平均之差的平方长度比较。
\solutionanswer
新条件平均\(M_Q=(2,2,4,8)\)，最小均方误差从\(5/2\)减为\([(-1)^2+1^2]/4=1/2\)，减少\(2\)。两种条件平均之差为\((0,0,-2,2)\)，平方长度\((4+4)/4=2\)，恰与误差减少量相同。


% ITEM: probability-exercise-12-03
\solutionitem{probability-exercise-12-03}{习题6.12（3） *}
若分割\(\mathcal Q\)比\(\mathcal P\)更细，用正交分解证明\(E(E(X\mid\mathcal Q)\mid\mathcal P)=E(X\mid\mathcal P)\)。
\solutionanswer
记两分组子空间为\(W_P\subseteq W_Q\)，相应投影为\(M_P,M_Q\)。\(X-M_Q\)正交于\(W_Q\)，从而正交于\(W_P\)；\(X-M_P\)也正交于\(W_P\)。相减得\(M_Q-M_P\perp W_P\)，而\(M_P\in W_P\)，所以它正是\(M_Q\)在\(W_P\)上的正交投影，即\(E(M_Q\mid\mathcal P)=M_P\)。
\par
